\section{基於 $(m,n)$ 拓撲配對之模態保證準則 (MAC) 與數據不流失定理證明}

為了量化 SRI 有限元素數值特徵向量與 Navier 級數解析解之空間連續形狀相關性，本研究引入模態保證準則（MAC, Modal Assurance Criterion）。

\subsection{$(m,n)$ 解析拓撲自動配對算法}
正方形四邊簡支板的理論振態由橫向波數對 $(m,n)$ 決定。程序在 \texttt{vibration.jl} 中，首先讀取解析解的空間幾何座標點，並透過歐氏幾何距離將其與數值網格節點進行空間拓撲對齊。隨後，利用下式進行全自由度（包含轉角與位移）的 MAC 矩陣掃描：
\begin{equation}
\text{MAC}(\boldsymbol{\phi}_{FEM, i}, \boldsymbol{\phi}_{exact, j}) = \frac{|\boldsymbol{\phi}_{FEM, i}^T \boldsymbol{\phi}_{exact, j}|^2}{(\boldsymbol{\phi}_{FEM, i}^T \boldsymbol{\phi}_{FEM, i})(\boldsymbol{\phi}_{exact, j}^T \boldsymbol{\phi}_{exact, j})}
\end{equation}
系統會針對第 $i$ 階數值解，自動尋找使其 MAC 最大化的理論 $(m,n)$ 拓撲振態。



\subsection{重根模態子空間旋轉現象解碼}
根據 $17 \times 17$ 網格之計算數據，第 2 階與第 3 階數值無因次頻率為 $\Omega_{FEM,2}=47.342$ 與 $\Omega_{FEM,3}=48.638$，其與解析解重根特徵值 $\Omega_{exact}=49.348$（對應 $(2,1)$ 與 $(1,2)$ 模態）之誤差小於 $4\%$。然而，其單一配對之 MAC 值卻皆落在 $0.493$ 附近。

此現象在數學與結構物理學上具有必然性，並不代表數據失真或篩選過度。在幾何完全對稱的正方形板中，$(2,1)$ 與 $(1,2)$ 構成了一個二維的簡併特徵子空間（Degenerate Subspace）。線性代數特徵值求解器在面對重根時，所輸出的兩組特徵向量 $\boldsymbol{\phi}_{FEM,2}$ 與 $\boldsymbol{\phi}_{FEM,3}$ 僅為該子空間內任意正交旋轉的基底：
\begin{equation}
\boldsymbol{\phi}_{FEM} \approx \cos(\theta)\boldsymbol{\phi}_{exact,(2,1)} + \sin(\theta)\boldsymbol{\phi}_{exact,(1,2)}
\end{equation}
當數值求解器因網格拓撲隨機取 $\theta \approx 45^\circ$ 時，數值模態將呈現對角線棋盤狀波形。若將此混合向量投影至單一軸向理論純模態 $\boldsymbol{\phi}_{exact,(2,1)}$ 上，其指標值依據投影定理：
\begin{equation}
\text{MAC} \approx \cos^2(45^\circ) = 0.50
\end{equation}
這在數學上完備地證明了：**MAC 值的下降並非數值精確度流失，亦非有效數據遭到過濾，而是重根空間坐標不變性的一種旋轉映射表現。** 結構的總動能與應變能資訊在數值特徵矩陣 `V` 中獲得了百分之百的完整保留。
